\section{Interpretation and further questions}\label{sec:discussion}

The local and companion results describe two different sources of sensitivity.
The normalized Jacobian law is determined by a first-order Gaussian calculation
and the determinant weight at a root. Its inverse moments have threshold four.
The companion overlap additionally contains the coefficient norm of a Lagrange
polynomial. For a fixed degree $n\ge2$ its moment threshold is again four, but degree
uniformity has threshold two. The radial estimates and the exact second moment
show why concentration of roots near the unit circle does not permit discarding
the complementary radii from conditioning averages.

This distinction also explains the relation to the limiting spectral measures
in \cite{barbarino2023monic,barbarino2022limit}. At a simple eigenvalue $\lambda$
of a matrix $M$, the spectral projector satisfies
\[
 \Pi_\lambda=\lim_{z\to\lambda}(z-\lambda)(zI-M)^{-1}
 =\frac{rw^*}{w^*r},\qquad \|\Pi_\lambda\|_2
 =\frac{\|r\|_2\|w\|_2}{|w^*r|}.
\]
Consequently an overlap estimate requires control of a resolvent residue.
Barbarino--Noferini's estimates for shifted monic companions
\cite[Lemma~A.11]{barbarino2023monic} concern fixed spectral parameters away
from the unit circle; they cannot simply be evaluated at a random eigenvalue,
where the shifted matrix is singular. Their decomposition into a deterministic
block shift and a Gaussian correction, followed by the Woodbury identity,
nevertheless suggests a useful route for studying such residues in block
companion models.

\subsection{A finite-degree illustration}

Table~\ref{tab:second-moments} evaluates the exact integral
\eqref{eq:comp-second-integral}. The conditional second moment on the unit
circle stays below two, while the unconditional second moment grows. These
values illustrate the proved formulas; they are not used to infer either
asymptotic or a uniform bound.

\begin{table}[htbp]
\centering
% Generated by research/paper-ac/second_moment_table.py
\begin{tabular}{rrrr}
\toprule
$n$ & $\Phi_n(1)$ & $\Estar\Lambda^2$ & $\Estar\Lambda^2-2\log n$\\
\midrule
1 & 1.000000 & 1.000000 & 1.000000 \\
2 & 1.333333 & 1.472800 & 0.086505 \\
8 & 1.777778 & 2.849863 & -1.309020 \\
32 & 1.939394 & 4.876708 & -2.054764 \\
128 & 1.984496 & 7.342059 & -2.362002 \\
\bottomrule
\end{tabular}

\caption{Deterministic quadrature of the exact second-moment formula.
Positive-coefficient Horner evaluation avoids cancellation at $r=1$.
Computations at 30 and 50 decimal digits agreed to more than 25 decimal places;
the table reports six. Degrees one and two also agree with their exact values
in Example~\ref{ex:comp-small}.}
\label{tab:second-moments}
\end{table}

The calculation is reproducible from
\texttt{research/paper-ac/second\_moment\_table.py}. A separate symbolic script
checks the cofactor polynomial, reversal, the conditional variance, and the
closed second-moment integrand in small degrees. These computations are
consistency checks of the identities, rather than replacements for the proofs.

\subsection{What a conditioning comparison establishes}

Corollary~\ref{cor:comp-relative} says that the geometric mean of the relative
companion-to-coefficient condition ratio is comparable to $n$. The ratio uses
the same scalar output but two different input spaces. It is compatible with
bounded average logarithmic eigenvector overlap because the Euclidean
coefficient-to-root condition has its own degree scale. Neither statement
alone bounds floating-point error in a complete algorithm.

An analysis of a particular implementation would have to control formation of
the companion matrix, the structure of the backward error from the eigensolver,
and any subsequent recovery. Equation~\eqref{eq:setup-composition} explains
the role of structure: a perturbation induced exactly by polynomial
coefficients has the original root derivative, whereas an arbitrary matrix
perturbation need not arise from a comparably small coefficient perturbation.
Other linearizations can change this comparison
\cite{toh1994pseudozeros,deteran2017eigenvalue}.

The selected-root tail also gives, by a union bound,
\[
 \Prob\{\max_{p(z)=0}\Lambda(p,z)>t\}\le\min\{1,Cn/t^2\}.
\]
The factor $n$ is part of this bound. Controlling a typical root, an expected
fraction of roots, and the worst root are different probabilistic objectives.

\subsection{The multivariate problem left by the local theory}

For systems, the representation question is now isolated in
$\Theta_{n,d}^{\B}$ from \eqref{eq:box-decomposition}. The local law and the
box metric estimates already hold in every dimension. What remains is to
control interpolation angles for an explicitly specified quotient basis or
algebraic construction. The univariate division map $D_z$ has a particularly
simple tridiagonal Gram matrix; a multivariate replacement must account for
relations between polynomial multiples and for basis selection.

Two extensions require additional ideas. First, bounds for a fixed bivariate
construction would begin to connect the local theory with multivariate
algebraic rootfinding. Any dependence on the dimension must be retained if
the objective is to compare with dimension-exponential worst cases. Second,
an additive perturbation model $P=P_0+\sigma Q$, where $P_0$ is fixed and $Q$
is a centered Gaussian system, produces noncentral conditional
derivatives and a noncentral determinant weight. The centered independence
law of Theorem~\ref{thm:local-universal} cannot be transferred to that smoothed
model without further analysis. The results proved here provide the centered
benchmark for these questions. Smoothed local conditioning and homotopy
complexity are already studied in the homogeneous Bombieri--Weyl setting
\cite{burgisser2011smale,armentano2016condition}; see also
Section~\ref{sec:local-classical}. The remaining task here is to control the
specified coefficient and representation conditions for the box ensemble
and its noncentered perturbations.

The noncentral Gaussian estimate of
\cite[Theorem~A.5]{barbarino2023monic} is a possible ingredient in this task,
not a substitute for the determinant-weighted law at a selected root.
Moreover, matrices produced by multivariate elimination generally reuse the
same input coefficients in several entries. Applying random matrix-polynomial
results to them requires treating these dependencies explicitly.
